% ============================================================
% anexos.tex — Anexos técnicos
% ============================================================

\section{Asignación de pines del ESP32-C3 SuperMini}
\label{anx:pines}

La Tabla~\ref{tab:pines} detalla la asignación de pines GPIO del ESP32-C3 SuperMini
utilizada en el prototipo NodeAlert IoT.

\begin{table}[htbp]
  \caption{\textit{Asignación de pines GPIO del ESP32-C3 SuperMini en el prototipo NodeAlert IoT.}}
  \label{tab:pines}
  \centering
  \small
  \begin{tabular}{@{}cll@{}}
    \toprule
    \textbf{GPIO} & \textbf{Función} & \textbf{Notas} \\
    \midrule
    0 & Relé canal 1 (salida) & Activo en nivel bajo \\
    1 & LED de presencia (salida) & Activo en nivel alto \\
    4 & I\textsuperscript{2}C SDA & Pull-up interno habilitado \\
    5 & I\textsuperscript{2}C SCL & Pull-up interno habilitado \\
    6 & LD2420 TX (UART1) & Transmisión hacia el radar \\
    7 & LD2420 RX (UART1) & Recepción desde el radar \\
    \bottomrule
  \end{tabular}
  \vspace{4pt}

  \noindent\textit{Nota.} Los pines de Wi-Fi y USB no se listan por ser internos al SoC
  y no requerir conexión externa. El bus I\textsuperscript{2}C opera a 100\,kHz.
\end{table}

\section{Esquema de tópicos MQTT y formato JSON}
\label{anx:mqtt}

El sistema utiliza la convención de nomenclatura de tópicos
\texttt{nodealert/<device\_id>/<tipo>}. El campo \texttt{device\_id} en el cuerpo
del mensaje JSON permite al servidor verificar la consistencia del origen del mensaje.

Los campos del mensaje de telemetría son opcionales a excepción de \texttt{device\_id}:
los campos de sensores en falla se omiten del JSON publicado. Esta decisión reduce
el tráfico de red y facilita la distinción entre «sensor ausente» y «sensor con
valor cero».

\section{Estructura del repositorio}
\label{anx:estructura}

\begin{lstlisting}[caption={Estructura de directorios del proyecto NodeAlert IoT.}, label={lst:estructura}]
nodered/
|-- README.md
|-- firmware/
|   `-- control-luz/           # Firmware ESP32-C3
|       |-- main/
|       |   |-- main.c         # Punto de entrada
|       |   |-- regla_luz.h    # Logica de control (funcion pura)
|       |   |-- tiempo.h       # Utilidades de tiempo
|       |   |-- config_local.h # Credenciales locales (no versionado)
|       |   `-- drivers/
|       |       |-- bh1750.c/.h
|       |       |-- sht31.c/.h
|       |       `-- ld2420.c/.h
|       |-- CMakeLists.txt
|       `-- sdkconfig           # Configuracion ESP-IDF
`-- raspberry/                  # Stack servidor
    |-- docker-compose.yml
    |-- generate-flows.js       # Generador de flujos Node-RED
    |-- mqtt/                   # Contenedor Mosquitto
    |   |-- Dockerfile
    |   |-- mosquitto.conf
    |   `-- aclfile
    |-- nodered/                # Contenedor Node-RED
    |   |-- Dockerfile
    |   `-- data/
    |       `-- flows.json
    `-- frontend/               # Contenedor React/Nginx
        |-- Dockerfile
        |-- nginx.conf
        |-- vite.config.js
        `-- src/
            |-- App.jsx
            |-- api.js
            `-- pages/
                |-- Dashboard.jsx
                |-- Nodes.jsx
                `-- NodeDetail.jsx
\end{lstlisting}

\section{Instrucciones de compilación del firmware}
\label{anx:compilacion}

Para compilar y flashear el firmware se requiere el ESP-IDF v6.1.0 correctamente
instalado. Los comandos de compilación son:

\begin{lstlisting}[language=bash, caption={Comandos de compilación y flasheo del firmware.}, label={lst:build}]
cd firmware/control-luz
# Configurar credenciales locales
cp main/config_local.ejemplo.h main/config_local.h
# Editar main/config_local.h con SSID, password y URI MQTT

# Compilar
idf.py set-target esp32c3
idf.py build

# Flashear y monitorear
idf.py -p /dev/ttyUSBx flash monitor
\end{lstlisting}

Para desplegar el stack servidor:

\begin{lstlisting}[language=bash, caption={Despliegue del stack servidor en Raspberry Pi.}, label={lst:deploy}]
cd raspberry
docker compose up -d --build
# La interfaz web estara disponible en http://<IP-RPi>:3001
# Node-RED en http://<IP-RPi>:1880
\end{lstlisting}

